% Artículo VII. Fragmento de cuerpo: corriente y respuesta Cartan--Holst.
% RESULTADO_RECUPERADO: acción, corriente variacional, inversa de Holst,
% torsión algebraica, respuesta cuadrática y normalización métrica.
% Fuentes leídas íntegramente:
% [II10] output/ARTICULO_II_REV10_EDICION_INTEGRADA_20260910/manuscrito/
%        sections/10_gravedad.tex, líneas 109--155.
% [G11c] integral activo 2249, fuente/manuscrito/sections/md/
%        11c_gravedad_holonomia_rev6.tex, líneas 214--334.
% [B028] integral activo, fuente/manuscrito/integracion_83/parte_iv/body/
%        B028_ch68_realizacion_cartan_holst_barbero_7df2afd9f675_11_gravedad_cartan_holst.tex,
%        líneas 76--218; copia idéntica conservada en II10/fuentes/propietarios_II/G03/.
%        SHA-256 e1e1f3325399b3ef09c331ed8268e3883f7bcc561cd7a560cf80d1be7eb1e91c.
% CERTIFICADO_NUEVO de resultados recuperados: inversión explícita del mapa
% torsión -> D(e wedge e), prueba de unicidad de contorsión y eliminación
% homogénea del término cuadrático. Las dos inversas se comprobaron sobre
% las 24 bases correspondientes mediante aritmética racional exacta.
% Precisión de dominio: la respuesta lineal en corriente y su eliminación
% cuadrática se enuncian para materia afín en la contorsión. Se conserva
% la dependencia del coeficiente respecto de la acción material concreta.
%
% DEPENDENCIAS DE INTEGRACIÓN (no constituyen remisiones bibliográficas):
% 1. vii:sec:realizacion-cartan: cotetrada, conexión y realización del
%    transporte enriquecido; incluir su prueba en el cuerpo precedente.
% 2. vii:sec:acoplamientos: genealogía de gamma_HMT, G_int, c_int,
%    homogeneidad dimensional y elección de la carta, con pruebas.
% 3. Para evaluar una especie material: acción S_m concreta, representación
%    de sus campos y cálculo de su corriente sigma. Este fragmento recibe
%    esa acción declarada y demuestra su respuesta variacional.
% 4. La identificación de una corriente de pantalla con sigma, su
%    contracción física y su normalización material tienen residencia propia.
%    Aquí no se introduce ni se selecciona una amplitud s_0.
% 5. El balance métrico utiliza la covariancia de la acción material y
%    sus ecuaciones de movimiento; no supone conservación separada del
%    tensor de espín y del tensor material de Levi--Civita.

\section{Courant de spin et réponse algébrique de Cartan--Holst}
\label{vii:sec:corriente-respuesta-cartan}

La réalisation du transport enrichi de la
section~\ref{vii:sec:realizacion-cartan} fournit la cotétrade et la
connexion ; les sections dimensionnelles et le paramètre d'incidence des
sections~\ref{v:ant:sec:barbero} et~\ref{v:ant:sec:gravity} fixent leurs couplages. Leur composition
détermine un problème variationnel du premier ordre. Le courant matériel
y occupe une position précise : il est la dérivée de l'action par rapport à
la connexion et détermine la réponse torsionnelle avec la même normalisation.

\subsection{Domaine géométrique et normalisation du courant}
\label{vii:subsec:dominio-corriente-cartan}

Nous travaillons sur une variété lisse orientée \(M\) de dimension quatre, munie d'un fibré lorentzien interne \(E\), d'une métrique \(\eta=\operatorname{diag}(-1,1,1,1)\) et d'une cotétrade non dégénérée \(e\in\Omega^1(M;E)\). Dans un repère local, nous écrivons \(e^I\), avec \(I=0,1,2,3\). La connexion métrique \(\omega^{IJ}=-\omega^{JI}\) agit au moyen de \(D_\omega\). Lorsque les champs \(\Psi\) sont spinoriels, une structure de spin et la représentation matérielle correspondante sont conservées. Nous définissons
\begin{equation}
 \Sigma^{IJ}=e^I\wedge e^J,\qquad
 F^{IJ}=\mathrm d\omega^{IJ}+\omega^I{}_{K}\wedge\omega^{KJ},
 \qquad T^I=D_\omega e^I.
 \label{vii:eq:campos-cartan}
\end{equation}
L'opérateur \(\star\) agit sur les indices bivectoriels internes,
\begin{equation}
 (\star X)^{IJ}=\tfrac12\epsilon^{IJ}{}_{KL}X^{KL},
 \qquad \star^2=-I,\qquad
 \mathsf P_\gamma=\star+\gamma^{-1}I.
 \label{vii:eq:operador-holst}
\end{equation}
On distingue ainsi cette dualité de la dualité extérieure des formes sur
\(M\). Les indices internes sont relevés et contractés avec \(\eta\).

Nous fixons un système de coordonnées avec \(G_{\rm int}>0\), \(c_{\rm int}>0\),
\(\gamma=\gamma_{\rm HMT}\in\mathbb R\setminus\{0\}\) et
\(\Lambda_{\rm int}\) constants sur celui-ci. Nous abrégeons
\(\kappa=8\pi G_{\rm int}/c_{\rm int}^4\) et \(c=c_{\rm int}\).
Ces sections restent fixes pendant la variation. L’action est
\begin{equation}
 \begin{split}
 S[e,\omega,\Psi]={}&
 \frac1{2\kappa c}\int_M\Sigma_{IJ}\wedge
                         (\mathsf P_\gamma F)^{IJ}
 -\frac{\Lambda_{\rm int}}{\kappa c}
                         \int_M\operatorname{vol}_e
 +S_{\rm m}[e,\omega,\Psi].
 \end{split}
 \label{vii:eq:accion-cartan-holst}
\end{equation}
La convention de volume et d'orientation est celle qui donne
\(\Sigma_{IJ}\wedge(\star F)^{IJ}=R\operatorname{vol}_e\)
dans le secteur de Levi--Civita. Le type dimensionnel préalable fait de chaque
terme une section d'action.

Pour la variation de la connexion, on utilise des variations à support compact à l'intérieur, ou on fixe \(\delta\omega|_{\partial M}=0\). On peut également employer une complétion de bord dont le problème variationnel annule le même terme de frontière. À cotétrade et champs matériels fixes, le courant de spin est la forme de degré trois, à valeurs dans les bivecteurs internes duaux, déterminée par
\begin{equation}
 \delta_\omega S_{\rm m}
 =-\tfrac12\int_M\delta\omega^{IJ}\wedge\sigma_{IJ}.
 \label{vii:eq:corriente-variacional}
\end{equation}
L'appariement avec une variation arbitraire à support compact fixe
\(\sigma\) de manière unique. Cette définition conserve l'unité
d'action du courant et le facteur de normalisation qui intervient dans
sa réponse géométrique.

\subsection{Variation de la connexion et inversion de Holst}
\label{vii:subsec:variacion-holst}

\begin{theorem}[Équation variationnelle de Cartan--Holst]
\label{vii:thm:ecuacion-cartan-holst}
Dans le domaine précédent, la stationnarité par rapport à la connexion
équivaut à
\begin{equation}
 D_\omega(\mathsf P_\gamma\Sigma)^{IJ}
       =\kappa c\,\sigma^{IJ}.
 \label{vii:eq:ecuacion-cartan-holst}
\end{equation}
\end{theorem}
\begin{proof}
La variation de la courbure est
\(\delta F=D_\omega\delta\omega\). La dualité interne est parallèle
et autoadjointe pour la contraction bivectorielle ; comme \(\gamma\) est
constant, \(D_\omega\mathsf P_\gamma=\mathsf P_\gamma D_\omega\).
En posant \(B_{IJ}=(\mathsf P_\gamma\Sigma)_{IJ}\), on obtient
\[
 B_{IJ}\wedge D_\omega\delta\omega^{IJ}
 =\mathrm d(B_{IJ}\wedge\delta\omega^{IJ})
   -D_\omega B_{IJ}\wedge\delta\omega^{IJ}.
\]
La dernière forme change de signe lorsqu'on permute ses facteurs de degrés trois
et un. Les conditions au bord éliminent l'intégrale du terme
exact. Le terme cosmologique est indépendant de \(\omega\).
Par conséquent,
\[
 \delta_\omega S=
 \frac1{2\kappa c}\int_M\delta\omega^{IJ}\wedge
       \bigl[D_\omega B_{IJ}-\kappa c\,\sigma_{IJ}\bigr].
\]
Le caractère arbitraire de la variation prouve l'équation, y compris son signe et son coefficient.
\end{proof}

\begin{lemma}[Inverse réel de l'opérateur de Holst]
\label{vii:lem:inversa-holst}
Pour \(\gamma\in\mathbb R\setminus\{0\}\),
\begin{equation}
 \mathsf P_\gamma^{-1}
 =\frac{\gamma^2}{1+\gamma^2}
                      (-\star+\gamma^{-1}I).
 \label{vii:eq:inversa-holst}
\end{equation}
\end{lemma}
\begin{proof}
Les deux facteurs sont des polynômes en \(\star\) et commutent. Leur produit
avant normalisation est
\[
 (\star+\gamma^{-1}I)(-\star+\gamma^{-1}I)
 =-\star^2+\gamma^{-2}I=(1+\gamma^{-2})I.
\]
Le dénominateur \(1+\gamma^2\) est positif dans le domaine réel fixé. La multiplication par \(\gamma^2/(1+\gamma^2)\) donne l'identité.
\end{proof}

Le courant détermine ainsi la forme bivectorielle de degré trois
\begin{equation}
 Y^{IJ}:=\kappa c\,
   \frac{\gamma^2}{1+\gamma^2}
   (-\star+\gamma^{-1}I)\sigma^{IJ},
 \qquad D_\omega\Sigma^{IJ}=Y^{IJ}.
 \label{vii:eq:corriente-dualizada}
\end{equation}

\subsection{Reconstruction ponctuelle de la torsion et de la contorsion}
\label{vii:subsec:reconstruccion-torsion}

La règle de Leibniz donne
\begin{equation}
 D_\omega\Sigma^{IJ}
 =T^I\wedge e^J-e^I\wedge T^J
 =:\mathsf L_e(T)^{IJ}.
 \label{vii:eq:mapa-torsion}
\end{equation}
En chaque point, \(\mathsf L_e\) agit de \(\Lambda^2T^*M\otimes E\) dans \(\Lambda^3T^*M\otimes\Lambda^2E\). Les deux espaces sont de dimension vingt-quatre. Soit \(E_I\) le repère dual de la cotétrade, de sorte que \(\iota_{E_I}e^J=\delta_I^J\).

\begin{theorem}[Inversion algébrique de l'équation torsionnelle]
\label{vii:thm:inversion-torsion}
Pour une cotétrade non dégénérée, \(\mathsf L_e\) est un isomorphisme.
Son inverse sur la forme \(Y\) s'écrit
\begin{equation}
 B^I=\sum_J\iota_{E_J}Y^{IJ},\qquad
 b=\tfrac14\sum_I\iota_{E_I}B^I,\qquad
 T^I=B^I-e^I\wedge b.
 \label{vii:eq:torsion-explicita}
\end{equation}
L'équation~\eqref{vii:eq:ecuacion-cartan-holst} détermine donc
la torsion ponctuellement à partir de \(e\) et du courant qui y
figure.
\end{theorem}
\begin{proof}
Soit d'abord \(Y=\mathsf L_e(T)\) et posons
\(t=\sum_J\iota_{E_J}T^J\). La contraction de
\eqref{vii:eq:mapa-torsion}, en utilisant le fait que \(T^I\) est de degré deux
et qu'il existe quatre éléments dans la cotétrade, donne
\[
 \sum_J\iota_{E_J}(T^I\wedge e^J)=2T^I,
 \qquad
 \sum_J\iota_{E_J}(e^I\wedge T^J)=T^I-e^I\wedge t.
\]
D'où \(B^I=T^I+e^I\wedge t\). Une seconde contraction produit
\[
 \sum_I\iota_{E_I}B^I=t+3t=4t.
\]
On retrouve \(t=b\) puis \(T^I=B^I-e^I\wedge b\).
En particulier, le noyau de \(\mathsf L_e\) est nul. L'égalité
des dimensions prouve sa surjectivité et rend la formule valable
pour toute forme \(Y\) du codomaine.
\end{proof}

La connexion se décompose de manière unique comme
\begin{equation}
 \omega^{IJ}=\mathring\omega^{IJ}(e)+K^{IJ},\qquad
 K^{IJ}=-K^{JI},\qquad T^I=K^I{}_{J}\wedge e^J,
 \label{vii:eq:contorsion}
\end{equation}
où \(\mathring\omega(e)\) est la connexion de Levi--Civita et \(K\) la contorsion. Pour expliciter l'inverse, nous écrivons \(K_{IJ}=K_{IJK}e^K\) et \(T_I=\tfrac12T_{IJK}e^J\wedge e^K\). Alors
\begin{equation}
 T_{IJK}=K_{IKJ}-K_{IJK},\qquad
 K_{IJK}=\tfrac12(T_{KIJ}-T_{IJK}-T_{JKI}).
 \label{vii:eq:contorsion-componentes}
\end{equation}
La seconde égalité s'obtient en combinant cycliquement la première et
en utilisant \(K_{IJK}=-K_{JIK}\) ; sa substitution restitue les composantes
de \(T\). La torsion et la contorsion sont donc deux
présentations équivalentes de la réponse algébrique à cotétrade fixée.

\begin{corollary}[Secteur sans courant et réponse linéaire minimale]
\label{vii:cor:desacoplamiento-torsional}
Si \(\sigma=0\), on obtient \(T=0\) et \(\omega=\mathring\omega(e)\). Pour un couplage matériel dont le courant \(\sigma[e,\Psi]\) est indépendant de \(K\), les équations~\eqref{vii:eq:corriente-dualizada}, \eqref{vii:eq:torsion-explicita} et \eqref{vii:eq:contorsion-componentes} fournissent une réponse unique \(K_*(e,\Psi)\), linéaire en \(\sigma\).
\end{corollary}
\begin{proof}
Les trois applications de reconstruction sont linéaires lorsque \(e\),
\(\gamma\) et \(\kappa c\) sont fixés. Le courant nul produit
\(Y=0\), puis \(T=0\) et enfin \(K=0\).
\end{proof}

Dans l'action minimale, l'équation de connexion contient cette réponse
ponctuelle. Son évolution s'obtient en faisant évoluer \(e\) et \(\Psi\) :
\(K_*=K_*(e,\Psi)\) change avec eux. Si une action matérielle conserve
une dépendance algébrique supplémentaire en \(K\), l'équation de Cartan
maintient sa forme variationnelle et doit être résolue avec ce courant ;
la réponse linéaire précédente correspond au domaine minimal spécifié.

\subsection{Élimination de la contorsion et contribution quadratique}
\label{vii:subsec:respuesta-cuadratica}

Considérons maintenant le couplage minimal affine en \(K\), exprimé par
\begin{equation}
 S_{\rm m}[e,\mathring\omega+K,\Psi]
 =S_{\rm m}[e,\mathring\omega,\Psi]
   -\tfrac12\int_M K^{IJ}\wedge\sigma_{IJ}[e,\Psi].
 \label{vii:eq:materia-afin-contorsion}
\end{equation}
Cette condition identifie l'action matérielle à laquelle s'applique
l'élimination quadratique ; le courant reste celui défini par
\eqref{vii:eq:corriente-variacional}.

Le développement de la courbure est
\begin{equation}
 F(\mathring\omega+K)
 =F(\mathring\omega)+D_{\mathring\omega}K+K\wedge K,
 \qquad (K\wedge K)^{IJ}=K^I{}_{L}\wedge K^{LJ}.
 \label{vii:eq:curvatura-contorsion}
\end{equation}
Comme \(D_{\mathring\omega}\Sigma=0\), le terme linéaire est la
dérivée extérieure de \(\Sigma_{IJ}\wedge(\mathsf P_\gamma K)^{IJ}\).
La contribution de bord correspondante reste explicite :
\begin{equation}
 S_\partial[e,K]=\frac1{2\kappa c}
     \int_{\partial M}\Sigma_{IJ}\wedge(\mathsf P_\gamma K)^{IJ}.
 \label{vii:eq:borde-contorsion}
\end{equation}
Dans un domaine intérieur, ou avec la complétion de bord compatible,
la densité d'action qui dépend de la contorsion est
\begin{equation}
 \mathcal Q_{e,\gamma}(K)-\tfrac12K^{IJ}\wedge\sigma_{IJ},
 \qquad
 \mathcal Q_{e,\gamma}(K)
  =\frac1{2\kappa c}\Sigma_{IJ}\wedge
                         \mathsf P_\gamma(K\wedge K)^{IJ}.
 \label{vii:eq:forma-cuadratica-contorsion}
\end{equation}
La densité \(\mathcal Q_{e,\gamma}\) est une forme différentielle de
degré quatre, homogène de degré deux en \(K\).

\begin{theorem}[Réponse effective quadratique en le courant]
\label{vii:thm:respuesta-cuadratica}
Sous \eqref{vii:eq:materia-afin-contorsion}, la substitution de la
solution algébrique \(K_*\) produit la densité effective de spin
\begin{equation}
 \mathcal L_{\rm spin}^{\rm eff}
 =-\tfrac14 K_*^{IJ}\wedge\sigma_{IJ}
 =-\mathcal Q_{e,\gamma}(K_*).
 \label{vii:eq:accion-efectiva-spin}
\end{equation}
À cotétrade et couplages fixés, elle est homogène de degré deux en
le courant. L'action effective conserve en outre
\(S_\partial[e,K_*]\), ou sa complétion de bord déclarée.
\end{theorem}
\begin{proof}
D'après le corollaire~\ref{vii:cor:desacoplamiento-torsional},
\(K_*\) est linéaire en \(\sigma\). L'équation de connexion est
précisément la stationnarité de
\eqref{vii:eq:forma-cuadratica-contorsion}. Comme celle-ci est algébrique,
on applique ponctuellement la variation radiale \(\delta K=K_*\) ;
elle peut aussi être multipliée par une fonction à support compact.
L'identité d'Euler pour une forme quadratique donne
\[
 2\mathcal Q_{e,\gamma}(K_*)
             -\tfrac12K_*^{IJ}\wedge\sigma_{IJ}=0.
\]
La substitution dans la densité d'action prouve \eqref{vii:eq:accion-efectiva-spin}. Si \(\sigma\mapsto\lambda\sigma\), alors \(K_*\mapsto\lambda K_*\) et la densité effective est multipliée par \(\lambda^2\). Le développement \eqref{vii:eq:curvatura-contorsion} identifie séparément le terme de bord qui accompagne cette élimination.
\end{proof}

Le courant concret, sa contraction lorentzienne, la dualité et les
conventions de l'action matérielle déterminent le signe et le coefficient
de cette contribution. L'identité quadratique fournit la règle de
composition qui permet de les évaluer une fois ces données spécifiées.

\subsection{Source métrique effective et bilan covariant}
\label{vii:subsec:fuente-metrica-efectiva}

La normalisation de la source métrique est fixée dans la même droite d'action. Après élimination de \(K\), nous définissons
\begin{align}
 \delta_g S_{\rm m}[e,\mathring\omega,\Psi]
 &=-\tfrac12\int_M\operatorname{vol}_e\,
                   \mathcal T^{S,\rm LC}_{\mu\nu}\,\delta g^{\mu\nu},
 \\
 \delta_g S_{\rm spin}^{\rm eff}
 &=-\tfrac12\int_M\operatorname{vol}_e\,
                   \mathcal U^{S,\rm spin}_{\mu\nu}\,\delta g^{\mu\nu}.
 \label{vii:eq:fuentes-normalizadas-accion}
\end{align}
On considère ici des variations intérieures, avec les termes de bord
traités conformément au problème variationnel fixé. Pour le couplage
affine, la seconde source est la variation métrique de
\eqref{vii:eq:accion-efectiva-spin}. La variation inclut la dépendance
du courant matériel à l'égard de \(e\) et de \(\Psi\).

\begin{proposition}[Équation métrique et conservation de la source totale]
\label{vii:prop:fuente-metrica-balance}
Avec les conventions précédentes, les équations métriques sont
\begin{equation}
 G_{\mu\nu}[g]+\Lambda_{\rm int}g_{\mu\nu}
 =\kappa c\bigl(\mathcal T^{S,\rm LC}_{\mu\nu}
                    +\mathcal U^{S,\rm spin}_{\mu\nu}\bigr)
 =\kappa\bigl(T^{\rm phys,LC}_{\mu\nu}
                    +U^{\rm phys,spin}_{\mu\nu}\bigr),
 \label{vii:eq:ecuacion-metrica-efectiva}
\end{equation}
où \(T^{\rm phys,LC}=c\mathcal T^{S,\rm LC}\) et \(U^{\rm phys,spin}=c\mathcal U^{S,\rm spin}\) lorsque la cotétrade prend ses valeurs dans la droite de longueur. Toute solution satisfait
\begin{equation}
 \mathring\nabla^\mu
   (T^{\rm phys,LC}_{\mu\nu}+U^{\rm phys,spin}_{\mu\nu})=0.
 \label{vii:eq:balance-metrico-total}
\end{equation}
\end{proposition}
\begin{proof}
L'élimination de \(K\) conserve les équations des variables
restantes : dans la règle de dérivation en chaîne, le terme qui multiplie
\(\delta K_*\) s'annule par l'équation de connexion. Dans le secteur
de Levi--Civita, la première identité de Bianchi annule le terme
\(\Sigma_{IJ}\wedge F^{IJ}\), et la partie gravitationnelle s'écrit
\((2\kappa c)^{-1}\int_M(R-2\Lambda_{\rm int})\operatorname{vol}_e\).
Sa variation intérieure est
\[
 \frac1{2\kappa c}\int_M\operatorname{vol}_e\,
        (G_{\mu\nu}+\Lambda_{\rm int}g_{\mu\nu})\delta g^{\mu\nu}.
\]
La somme avec \eqref{vii:eq:fuentes-normalizadas-accion} donne \eqref{vii:eq:ecuacion-metrica-efectiva}. L'identité de Bianchi contractée pour \(\mathring\nabla\), la compatibilité métrique et la constance de \(\kappa\) et de \(\Lambda_{\rm int}\) dans la coordinatisation produisent \eqref{vii:eq:balance-metrico-total}.
\end{proof}

Avant d'éliminer la connexion, les identités géométriques conservent
la forme
\begin{equation}
 D_\omega T^I=F^I{}_{J}\wedge e^J,\qquad D_\omega F^{IJ}=0.
 \label{vii:eq:bianchi-cartan}
\end{equation}
La première est \(D_\omega^2e=F\wedge e\) ; la seconde s'obtient
en développant \(D_\omega(\mathrm d\omega+\omega\wedge\omega)\)
et en annulant les termes. Dans une action matérielle covariante, le
bilan métrique total coïncide avec l'identité de Noether sur les
équations matérielles. Géométrie et matière conservent ainsi une source
conjointe ; la répartition entre sa contribution de Levi--Civita et sa
contribution de spin est déterminée par l'action spécialisée.
